\scanpage{67}
- \mspage{108} -
\hfill n°371
\[
\langle z,u\llcorner(s\wedge t)\rangle=\langle(s\wedge t)\wedge z,u\rangle=\langle t\wedge z,u\llcorner s\rangle=\langle z,(u\llcorner s)\llcorner t\rangle
\]
en vertu de l'associativité de la multiplication dans $\bigwedge(M)$ ; d'où
\[
(31)\qquad u\llcorner(s\wedge t)=(u\llcorner s)\llcorner t.
\]
On peut donc dire que l'application $(t,u)\longrightarrow u\llcorner t$ définit (avec l'addition) une structure de $\bigwedge(M)$-\uline{module à droite gradué} sur $\bigwedge(M)^{*\mathrm{gr}}$ (quand on munit $\bigwedge(M)$ de la graduation \uline{opposée} à celle définie au n°1).

Soit maintenant $f:M\longrightarrow N$ une application $A$-linéaire ; pour $z,t$ dans $\bigwedge(M)$ et $u'\in\bigwedge(N)^{*\mathrm{gr}}$, on a par définition
\[
\begin{aligned}
\langle t\wedge z,{}^{\mathrm t}(\bigwedge(f))(u')\rangle&=\langle(\bigwedge(f))(t\wedge z),u'\rangle=\\
&=\langle(\bigwedge(f))(t)\wedge(\bigwedge(f))(z),u'\rangle=\langle(\bigwedge(f))(z),u'\llcorner(\bigwedge(f))(t)\rangle=\\
&=\langle z,{}^{\mathrm t}(\bigwedge(f))(u'\llcorner(\bigwedge(f))(t)\rangle
\end{aligned}
\]
d'où
\[
(32)\qquad{}^{\mathrm t}(\bigwedge(f))(u'\llcorner(\bigwedge(f))(t))={}^{\mathrm t}(\bigwedge(f))(u')\llcorner t.
\]
En particulier, si $f$ est un \uline{isomorphisme} de $M$ sur $N$, on a
\[
(33)\qquad{}^{\mathrm t}(\bigwedge(f))^{-1}(u\llcorner t)=({}^{\mathrm t}(\bigwedge(f))^{-1}(u))\llcorner((\bigwedge(f))(t))
\]
quels que soient $t\in\bigwedge(M)$, $u\in\bigwedge(M)^{*\mathrm{gr}}$ (formule qui résulte aussi du transport de structure par $f$).

* Nous verrons au §6, n°4 que si $u=x_1^*x_2^*\ldots x_q^*$, où les $x_i^*$ appartiennent à $M^*$, on a, pour tout $t\in\bigwedge^pM$ ($p\leq q$)
\[
(34)\qquad(x_1^*x_2^*\ldots x_q^*)\llcorner t=\sum_H\unsure{\rho|\varepsilon}_{H,K}\langle t,x_H^*\rangle x_K^*
\]
où $H$ parcourt l'ensemble des parties à $p$ éléments de $\lbrack\!\lbrack1,q\rbrack\!\rbrack$, $K$ est le complémentaire de $H$ dans $\lbrack\!\lbrack1,q\rbrack\!\rbrack$ et, pour toute partie $J$ de $\lbrack\!\lbrack1,q\rbrack\!\rbrack$, on pose $x_J^*=x_{i_1}^*x_{i_2}^*\ldots x_{i_m}^*$ si $(i_h)_{1\leq h\leq m}$ est la suite strictement croissante des éléments de $J$. *

Montrons maintenant que, pour $u\in\bigwedge(M)^{*\mathrm{gr}}$, il existe un endomorphisme canonique $s\longrightarrow u\lrcorner s$ du $A$-module $\bigwedge(M)$ tel que l'on ait
\margin{66}
